\documentclass[11pt]{article}
\usepackage[a4paper,margin=1in]{geometry}
\usepackage{amsmath,amssymb,amsthm,booktabs,microtype}
\usepackage[hidelinks]{hyperref}
\newtheorem{proposition}{Proposition}
\title{A Fixed-Source Cardinality-Monotonicity Obstruction\
for Native Prime-Shell Budgets}
\author{Liang Wang\\
School of Mathematics and Statistics\\
Huazhong University of Science and Technology (HUST), Wuhan, China}
\date{29 August 2026}
\begin{document}
\maketitle

\begin{abstract}
The source-first growing-grid audit of TPC-302 found a stable gap between a
signed weighted target and an all-positive control, but did not determine how
the weighted native budget itself changes with shell size.  We test the
stronger cardinality-only monotonicity shortcut on a fixed source spine:
$(N,H,z)=(512,58,5)$ and $Q=50,60,70,90$, with shell cardinalities
$10,13,15,17$.  We use the three TPC-302 relative RMS tolerances, two kernel
exponents, and three source normalizations.  Outward interval comparisons give
21 strict descents and 33 strict ascents among 54 adjacent transitions, with no
unresolved case.  All 18 parameter series are nonmonotone, and nine descents
occur at an unchanged common profile prefix.  The strongest same-prefix
contraction is below $0.284422$.  This is a scoped finite obstruction, not a
refutation of an eventual asymptotic lower bound.
\end{abstract}

\section{Question}

TPC-302 transported a native weighted/positive budget gap to a 34-row
growing/control grid and regenerated its signed target from the physical Gram
matrix \cite{tpc302}.  A tempting next shortcut would be to infer that the
weighted native budget must grow whenever shell cardinality grows.  Such an
inference is stronger than the TPC-302 gap statement and needs its own test.

We freeze source scale $N=512$, height $H=58$, comparison cutoff $z=5$, and
compare the moving shells indexed by $Q=50,60,70,90$.  Their cardinalities are
$10,13,15,17$.  The shells are not nested; the claim under attack is therefore
cardinality-only monotonicity, not monotonicity under set inclusion.

For a common profile prefix $k$ and target $b$, TPC-302 uses
\[
 B_{k,\tau}(b)=\min\{c^TM_kc:\|V_kc-b\|_2\leq\tau\|b\|_2\}.
\]
We read the common-prefix weighted budget, normalized by
\(\|\beta\|_2^2\), \(\operatorname{tr}(M_k)/k\), or $M_k[1,1]$.

\section{Interval certificate}

\begin{proposition}[Certified order]
If $B_j\in[L_j,U_j]$ and $B_{j+1}\in[L_{j+1},U_{j+1}]$, then
$U_{j+1}<L_j$ implies $B_{j+1}<B_j$, while $L_{j+1}>U_j$ implies
$B_{j+1}>B_j$.
\end{proposition}
\begin{proof}
Every represented value lies between its interval endpoints.  The strict
endpoint inequalities place the two values in disjoint ordered intervals.
\end{proof}

\begin{proposition}[Finite scoped refutation]
One certified descent refutes nondecreasing cardinality-only budget on the
declared finite Q-spine.
\end{proposition}
\begin{proof}
The definition of nondecreasing requires every adjacent pair to have the
later value at least the earlier value.  A certified strict descent violates
that condition.
\end{proof}

A descent with equal common prefix $k$ is especially informative: it cannot
be blamed on adding or removing profile directions.  It can still reflect the
moving physical shell or the source-first change in the signed target.

\section{Results}

\begin{table}[t]
\centering\small
\caption{Adjacent transition census on the fixed-source spine.}
\begin{tabular}{lrrr}
\toprule
quantity & count & total & fraction\\
\midrule
certified descents & 21 & 54 & $21/54$\\
certified ascents & 33 & 54 & $33/54$\\
unresolved overlaps & 0 & 54 & $0/54$\\
nonmonotone parameter series & 18 & 18 & $18/18$\\
same-prefix descents & 9 & 54 & $9/54$\\
\bottomrule
\end{tabular}
\end{table}

The strongest descent overall is the $Q=60\to70$ transition at exponent 2,
$\tau=1/2$, normalized by $\|\beta\|_2^2$; its right/left interval ratio is
below $0.224974$.  The strongest same-prefix descent is also $Q=60\to70$ at
exponent 2 and $\tau=3/4$, with both common prefixes equal to $k=6$ and ratio
below $0.284422$.  The three normalizers give the same qualitative result,
and each exponent--tolerance--normalizer series contains both an ascent and a
descent.

These are interval-certified finite comparisons inherited from a frozen
source-first certificate.  They do not say that a suitably normalized budget
cannot eventually obey a lower bound, nor do they turn moving-shell data into
a nested-shell theorem.

\section{Obstruction and next question}

The result closes a useful but invalid shortcut: shell cardinality alone does
not control the direction or size of the native weighted budget on this
declared spine.  The same-prefix cases further show that profile dimension is
not the only explanation.  The next natural experiment is to transport the
weighted sign labels across overlapping shells, separating target-label
switching from the physical operator/shell change.  Uniform profile-budget
growth, arithmetic $L^2$, fixed-power credit, full Gate B, and a twin-prime
conclusion remain open.

\section*{Reproducibility and scope}

The canonical TPC-303 certificate stores every interval and transition.  An
independent checker reconstructs the census from TPC-302 without importing
the producer, and a stress suite checks the interval logic.  The
Session-named Route-A/Route-B evaluator files are absent from this checkout;
no official evaluator pass is asserted.

\bibliographystyle{plain}
\bibliography{references}
\end{document}
